\documentclass[12pt]{article}
\usepackage[utf8]{inputenc}
\usepackage[T1]{fontenc}
\usepackage[margin=1in]{geometry}
\usepackage{setspace}
\usepackage{xcolor}
\usepackage[colorlinks=true, urlcolor=blue!60!black, linkcolor=blue!60!black]{hyperref}
\setlength{\parindent}{0em}
\setlength{\parskip}{0.6em}
\renewcommand{\baselinestretch}{1.0}
\begin{document}
\begin{center}
{\LARGE \textit{Moore v. United States}, 602 U.S. 572 (2024)\par}
\vspace{0.5em}
{\large Case Brief --- Two-Page Edition\par}
\end{center}

\section{Memory Jogger}
Congress may attribute an entity's realized, undistributed income to its shareholders or partners and tax them without apportionment; the Court upheld the 2017 Mandatory Repatriation Tax but reserved whether realization is constitutionally required (\href{https://supreme.justia.com/cases/federal/us/602/22-800/}{Justia}).

\section{Detailed Case Facts}
Charles and Kathleen Moore owned 13\% of KisanKraft, an Indian CFC that reinvested its profits. The 2017 MRT taxed 10\% U.S. shareholders on their shares of accumulated post-1986 foreign earnings; the Moores paid \$14,729.

\section{Procedural History}
Their refund suit was dismissed; the Ninth Circuit affirmed, saying ``realization of income is not a constitutional requirement.''

\section{Judicial Votes}
Affirmed, 7--2. Kavanaugh for the Court (Roberts, Sotomayor, Kagan, Jackson); Jackson concurring; Barrett, with Alito, concurring in the judgment; Thomas, with Gorsuch, dissenting.

\section{Holding}
The MRT taxes income realized by the CFC and attributed to its shareholders, which Congress may do without apportionment when the entity is not taxed on it. The Court reserved taxes on wealth, double taxation of entity and owners, and unrealized appreciation.

\section{Analysis of Opinions}
\textbf{Kavanaugh:} The Sixteenth Amendment restored income taxes to the unapportioned class; precedents and practice since 1864 permit taxing owners on entity income, and the Moores' theory would cause ``fiscal calamity.'' \textbf{Jackson:} the Amendment has no realization requirement. \textbf{Barrett:} realization is required and the Moores realized nothing, but they failed to distinguish Subpart F. \textbf{Thomas:} income must be realized by the taxpayer; the MRT taxed mere ownership.

\section{Examples: Future Applications}
\textbf{Same-Side:} Partners taxed on profits their partnership retains. \textbf{Opposite-Side:} A tax on billionaires' unrealized appreciation falls outside \textit{Moore}. \textbf{Fence-Sitter:} A pass-through tax on all public-company shareholders tests how far attribution reaches.

\section{Critique}
Critics say the Court ducked the question presented to avoid ``fiscal calamity,'' leaving a vague attribution standard; the split reflects formalism versus economic substance.

\section{Key Quotations}
\begin{quote}
``So the precise and narrow question that the Court addresses today is whether Congress may attribute an entity's realized and undistributed income to the entity's shareholders or partners, and then tax the shareholders or partners on their portions of that income. This Court's longstanding precedents, reflected in and reinforced by Congress's longstanding practice, establish that the answer is yes.''
\end{quote}
\noindent --- Justice Kavanaugh, Opinion of the Court.

\end{document}
